% titel: Blatt 6, A3: Lewis-Formel und Summenformel

\subsection*{a) Oktettregel}
Atome bilden so viele Bindungen, dass sie 8 Valenzelektronen erreichen
(Edelgaskonfiguration), H nur 2. Ausnahmen: $\mathrm{B}$ (Elektronenmangel, nur 6 Valenzelektronen, z.\ B.
$\mathrm{BF_3}$), $\mathrm{Be}$, Oktetterweiterung ab
Periode~3 ($\mathrm{PCl_5}$); als Radikal mit ungerader Elektronenzahl gilt
$\mathrm{NO}$.

\subsection*{b) Lewis-Struktur und Summenformel je Verbindung}
\begin{center}
\begin{tabular}{@{}l l l@{}}
  Verbindung & Summenformel & Lewis-Struktur-Merkmale \\
  \hline
  Wasser      & $\mathrm{H_2O}$ & O mit 2 H und 2 freien Paaren \\
  Ammoniak    & $\mathrm{NH_3}$ & N mit 3 H und 1 freien Paar \\
  Glycin      & $\mathrm{C_2H_5NO_2}$ &
     Carbonsäure- und Amin-Gruppe; N dreibindig, O als $\mathrm{OH}$ und als $\mathrm{=O}$ \\
  Lachgas     & $\mathrm{N_2O}$ & lineares $\mathrm{N=N=O}$ mit Resonanz \\
  Blausäure   & $\mathrm{HCN}$ & $\mathrm{H-C\equiv N}$ \\
  Formaldehyd & $\mathrm{CH_2O}$ & C mit $\mathrm{=O}$ und 2 H \\
  Salzsäure   & $\mathrm{HCl}$ & $\mathrm{H-Cl}$, ein freies Paar je Atom \\
  Ethanol     & $\mathrm{C_2H_5OH}$ & $\mathrm{CH_3-CH_2-OH}$ (O zweibindig) \\
  Aceton      & $\mathrm{C_3H_6O}$ & $\mathrm{CH_3-CO-CH_3}$ ($\mathrm{C=O}$) \\
\end{tabular}
\end{center}

(Auf dem Blatt ist für Ethanol irrtümlich $\mathrm{C_5H_5OH}$ geschrieben;
richtig ist $\mathrm{C_2H_5OH}$, die Summenformel $\mathrm{C_2H_6O}$.)

\begin{ergebnis}
Alle Verbindungen erfüllen die Oktettregel (je Atom 8, H 2); die Summenformel
zählt die Atome, die Lewis-Struktur zeigt Bindungen und freie Paare.
\end{ergebnis}

\begin{kontrolle}
Wertigkeitsprobe: $\mathrm{N}$ dreibindig + 1 freies Paar, $\mathrm{O}$
zweibindig, $\mathrm{C}$ vierbindig, H einbindig. Glycin
$\mathrm{C_2H_5NO_2}$ passt mit Carboxyl-$\mathrm{OH}$ und Amin-Gruppe zu dieser
Regel; kein Atom überbindet.
\end{kontrolle}

\begin{merksatz}
Oktettregel: 8 Elektronen je Atom (H 2). Summenformel zählt Atome;
Lewis-Struktur zeigt zusätzlich Bindungs- und freie Paare. Ausnahmen:
$\mathrm{B}$ (6), Oktetterweiterung ab Periode 3, Radikale.
\end{merksatz}
